\section{The box-product of cosimplicial objects}
The aim of this section is to introduce the box product $\square$ for cosimplicial objects in a monoidal category $(\mathsf{C}, \otimes, \mathbf{1})$ as first introduced by Batanin \cite{Bat-box}. For nice categories $\mathsf{C}$ (e.g.,  $\mathsf{C}$ closed, symmetric monoidal), the box product endows $\mathsf{C}^{\bdel}$ with a monoidal structure, and cosimplicial objects which admits a monoidal pairing with respect to $\square$ inherit an $A_{\infty}$-monoidal pairing on their totalizations (see, e.g.,  McClure-Smith \cite[3.1]{MS-box}).

Our main use of the box product will be to produce a homotopy-coherent (i.e., $A_{\infty}$-) composition on the derivatives of the identity in $\Cal{O}$-algebras, modeled as $\Tot C(\Cal{O})$, by demonstrating a natural pairing $C(\Cal{O})\square C(\Cal{O})\to C(\Cal{O})$ (Example \ref{ex:C(O)-box}).

\begin{definition} \label{def:box-product} 
Let $(\mathsf{C}, \otimes, \mathbf{1})$ be a monoidal category and $X,Y\in \textsf{C}^{\bdel}$. Define their \textit{box product} $X\square Y\in \textsf{C}^{\bdel}$ at level $n$ by \[
(X\square Y)^n:=
\colim \left( \xymatrix{ 
    \displaystyle{\coprod_{p+q=n}} X^{p}\otimes Y^{q} & 
    \ar@<-1 ex>[l] \ar@<.5 ex>[l] 
    \displaystyle{\coprod_{r+s=n-1}} X^{r}\otimes Y^{s} 
}\right) \]
where the maps are induced by $\id\otimes d^0$ and $d^{r+1}\otimes \id$. The object $X\square Y$ inherits cosimplicial structure via coface maps $d^i\colon (X\square Y)^{n}\to (X\square Y)^{n+1}$ induced by
\[\begin{cases} 
X^p\otimes Y^q \xrightarrow{\;\; d^i\otimes \id\;\;} X^{p+1}\otimes Y^{q} & i \leq p
\\
X^p\otimes Y^q \xrightarrow{\;\; \id \otimes d^{i-p} \;\;} X^{p}\otimes Y^{q+1} & i > p
\end{cases}
\]
and codegeneracy maps $s^j\colon (X\square Y)^n \to (X\square Y)^{n-1}$ induced by
\[\begin{cases} 
X^p\otimes Y^q \xrightarrow{\;\; s^j\otimes \id\;\;} X^{p-1}\otimes Y^{q} & j < p
\\
X^p\otimes Y^q \xrightarrow{\;\; \id \otimes s^{j-p} \;\;} X^{p}\otimes Y^{q-1} & j \geq p
\end{cases}
\]
see also Ching-Harper \cite[\textsection 4]{Ching-Harper-Koszul}.
\end{definition}

\begin{remark}Note, $(X\square Y)^0\cong X^0\otimes Y^0$, $(X\square Y)^{1}$ and $(X\square Y)^{2}$ may be computed as the colimits of \[ \xymatrix{ X^0 \otimes Y^1 & \\ X^0 \otimes Y^0 \ar[r]_-{d^1\otimes \id} \ar[u]^-{\id \otimes d^0} & X^1 \otimes Y^0} \; \raisebox{-8ex}{\text{ and }}  \; \xymatrix{ 
    X^0 \otimes Y^2 & &
    \\ 
    X^0 \otimes Y^1 \ar[r]_-{d^1\otimes \id} \ar[u]^-{\id \otimes d^0} 
    & 
    X^1 \otimes Y^1 & \\
    & X^1 \otimes Y^0 \ar[u]^-{\id \otimes d^0} \ar[r]_-{d^2\otimes \id} 
    &
    X^2 \otimes Y^0}
    \] respectively, and in general $(X\square Y)^n$ may be computed as the colimit of the staircase diagram \begin{equation}\xymatrix{
    X^0 \otimes Y^n & &\\
    X^0 \otimes Y^{n-1} \ar[r]_{d^1\otimes \id} \ar[u]^{\id\otimes d^0} & X^1 \otimes Y^{n-1} &\\
    & \ddots \ar[u] \ar[r]& X^{n-1}\otimes Y^1 &\\
    && X^{n-1}\otimes Y^0 \ar[u]^{\id\otimes d^0} \ar[r]_{d^n\otimes \id} &X^n\otimes Y^0 
}\end{equation}
\end{remark}

In particular, if $(\mathsf{C},\otimes, \mathbf{1})$ is closed, symmetric monoidal then $\square$ defines a monoidal category $(\textsf{C}^{\bdel}, \square, \underline{\mathbf{1}})$, here $\underline{\mathbf{1}}$ is the constant cosimplicial object on the unit $\mathbf{1}\in\mathsf{C}$ (see, e.g.,  Batanin \cite{Bat-box}). %We include the following as a motivating example for our $A_{\infty}$-operad structure on $\partial_* \Id$. 

\begin{example} \label{ex:C(O)-box}
Recall the cosimplicial symmetric sequence $C(\Cal{O})=J^{(\bullet+1)}$ from \eqref{al:CO}. We observe that $C(\Cal{O})$ admits a pairing $C(\Cal{O}) \boxcirc C(\Cal{O}) \xrightarrow{m} C(\Cal{O})$, where $\boxcirc$ denotes the box product in $\SymSeq^{\bdel}_{\Spt}$, induced as follows. Let $c$ denote the operad composition map $c\colon J\circ J\to J$. Then, \begin{align*} 
    &(C(\Cal{O}) \boxcirc C(\Cal{O}))^0\cong J \circ J \xrightarrow{c} J=C(\Cal{O})^0 \end{align*}
    
For level $1$ we observe that there are maps \[ m_{0,1} \colon J\circ J\circ_{\Cal{O}} J \to J\circ_{\Cal{O}} J \; \text{ and } \; m_{1,0} \colon J \circ_{\Cal{O}} J \circ J \to J\circ_{\Cal{O}} J \]
induced by $J\circ J\to J$ which induces $m$ via the following commuting square
\begin{equation}\xymatrix{
    J \circ J\circ_{\Cal{O}} J \ar[r]^-{m_{0,1}}
    & 
    J \circ_{\Cal{O}} J 
    \\
    J \circ J \ar[u]^-{\id\circ  d^0} \ar[r]^-{d^1\circ \id}
    &
    J\circ_{\Cal{O}} J\circ J \ar[u]^-{m_{1,0}}
} 
\end{equation}

More generally, there are maps of the form \[m_{p,q}\colon J^{(p)} \circ J^{(q)} \to J^{(p+q)} \;\; \text{for $p+q=n$, $p,q\geq 0$}\] induced by $c$, which induces the pairing $m$ at level $n$. 
\end{example}

\begin{remark}
    The above construction is entirely analogous to the following example found in McClure-Smith \cite{MS-box} that the based loop space $\Omega X$ of $X\in \Top$ admits an $A_{\infty}$ composition induced by an underlying $\square$-pairing. In this case, $\Omega X$ is modeled as the totalization of the cobar complex $c(X)$ built with respect to the natural comultiplication (with coaugmentation) given by the diagonal $\delta\colon X\to X\times X$. 
    
    It follows that $c(X)^{p}\cong X^{\times p}$ and the pairing $c(X)\square c(X)\to c(X)$ is induced by the natural isomorphisms $X^{\times p}\times X^{\times q}\cong X^{\times p+q}$. Further, McClure-Smith  show that $\Tot c(X)$ is an algebra over the (nonsymmetric) coendmorphism operad on $\Delta^{\bullet}$, i.e., \[ \mathbb{A}[n]=\Map_{\Dres}^{\mathsf{Top}}\left(\Delta^{\bullet}, (\Delta^{\bullet})^{\square n} \right)\] which satisfies $\mathbb{A}[0]=\pt$ and $\mathbb{A}[n]\xrightarrow{\sim} \pt$ for $n\geq 1$ (in fact $\Delta^n$ and $(\Delta^{\bullet} \square \Delta^{\bullet})^n$ are \textit{homeomorphic}), and that with respect to this structure $\Tot c(X)\simeq \Omega X$ as $A_{\infty}$-monoids.  \label{ex:loop-space-A-infty-monoid}
\end{remark}


\section{The box product in $\SymSeq^{\bdel}$}
Our aim now is to build a framework in which we can work with the structure captured by Example \ref{ex:C(O)-box}, e.g.,  by considering the box-product in the category of cosimplicial objects in $(\SymSeq_{\mathsf{C}}, \circ, I)$ of symmetric sequences for $(\mathsf{C},\otimes, \mathbf{1})$ some closed symmetric monoidal category. 

The main difficulty is that the composition product of symmetric sequences does not always commute with colimits taken in the right hand entry. That is, for $B\colon \mathcal{I}\to \SymSeq_{\mathsf{C}}$ a small diagram and $A\in \SymSeq_{\mathsf{C}}$, the universal map \begin{equation} \label{eq:colim-circ}\colim_{i\in \mathcal{I}} (A\circ B_i) \to A \circ (\colim_{i\in \mathcal{I}} B_i) \end{equation} is not an isomorphism in general. Thus the box-product fails to be strictly monoidal in this setting. 

Let us write $\SymSeq=\SymSeq_{\mathsf{C}}$ and $\boxcirc$ for the box-product in $\SymSeq^{\bdel}$ (in words we refer to $\boxcirc$ as the \textit{box-circle product}). Let $\Cal{X},\Cal{Y},\Cal{Z}, \dots$ be cosimplicial symmetric sequences. We will systematically interpret expressions of the form $\Cal{X} \boxcirc \Cal{Y} \boxcirc \Cal{Z}$ to be expanded \textit{from the left}, i.e., 
    \[\Cal{X} \boxcirc \Cal{Y} \boxcirc \Cal{Z} := (\Cal{X} \boxcirc \Cal{Y}) \boxcirc \Cal{Z}, \; \Cal{X} \boxcirc \Cal{Y} \boxcirc \Cal{Z} \boxcirc \Cal{W} := ((\Cal{X} \boxcirc \Cal{Y}) \boxcirc \Cal{Z} ) \boxcirc \Cal{W}, \;\dots \]
and note that via the universal map in (\ref{eq:colim-circ}) there is always a canonical comparison map $\theta$ of the form    
    \begin{equation} \label{eq:theta} \theta\colon  \Cal{X}\boxcirc \Cal{Y}\boxcirc \Cal{Z}=(\Cal{X} \boxcirc \Cal{Y}) \boxcirc \Cal{Z}  \to \Cal{X} \boxcirc (\Cal{Y} \boxcirc \Cal{Z})\end{equation}
which likely fails to be invertible. However, $\theta$ is sufficient to provide a suitable description of monoids with respect to $\boxcirc$, i.e., Definition \ref{boxcirc-monoid}, below. First, we note that the unit $I\in \SymSeq_{\mathsf{C}}$ induces a unit $\underline{I} \in \SymSeq_{\mathsf{C}}$ as the constant cosimplicial object on $I$ in that there are isomorphisms \[ \Cal{X} \boxcirc \underline{I} \cong \Cal{X} \cong \underline{I} \boxcirc \Cal{X}.\] 

For instance, the right isomorphism is obtained by noting that for any $p,q$ the map $d^{p+1}\circ\id$ in the following \[ \xymatrix{
        \underline{I}^p \circ \Cal{X}^{q+1} 
        & \\
        \underline{I}^p \circ \Cal{X}^q \ar[u]^-{\id\circ d^0} \ar[r]^-{d^{p+1}\circ \id}
        &
        \underline{I}^{p+1} \circ \Cal{X}^{q} 
        } \]
    is just the identity (and hence has an inverse). Therefore, the inclusion of the vertex $\underline{I}^0\circ \Cal{X}^n$ into the diagram defining $(\Cal{X}\boxcirc \Cal{X})^n$ is right cofinal (i.e., induces an isomorphism on colimits). 
    

\begin{definition} \label{boxcirc-monoid} By $\boxcirc$-monoid in $\SymSeq^{\bdel}$, we mean a cosimplicial symmetric sequence $\Cal{X}$ together with maps $m\colon \Cal{X} \boxcirc \Cal{X} \to \Cal{X}$ and $u\colon \underline{I} \to \Cal{X}$ so that the following associativity (\ref{eq:assoc-boxcirc}) and unitality (\ref{eq:unit-boxcirc}) diagrams commute
    \begin{equation} \label{eq:assoc-boxcirc} \xymatrix{  
        \Cal{X} \boxcirc \Cal{X} \boxcirc \Cal{X} \ar[r]^{\theta} \ar[d]^-{=}&
        \Cal{X} \boxcirc (\Cal{X} \boxcirc \Cal{X}) \ar[r]^-{\id \boxcirc m} &
        \Cal{X} \boxcirc \Cal{X} \ar[d]^-{m} \\
        (\Cal{X} \boxcirc \Cal{X}) \boxcirc \Cal{X} \ar[r]^-{m \boxcirc \id} &
        \Cal{X} \boxcirc \Cal{X} \ar[r]^-{m} &
        \Cal{X}
        }\end{equation}
    and
   \begin{equation} \label{eq:unit-boxcirc} \xymatrix{ 
        \Cal{X} \boxcirc \underline{I} \ar[r]^-{\id\boxcirc u} \ar[dr]_-{\cong}
        &
        \Cal{X} \boxcirc \Cal{X} \ar[d]^-{m}
        & 
        \underline{I}\boxcirc \Cal{X} \ar[l]_-{u\boxcirc \id} \ar[dl]^-{\cong}\\
        & \Cal{X}  &
        } 
    \end{equation}
\end{definition} 

\begin{remark}
We remark that in the language of Ching \cite{Ching-oplax} (see also Day-Street \cite{Day-St}), $\boxcirc$ admits a \textit{normal oplax monoidal structure} by defining \[ \Cal{X}_1\boxcirc \cdots \boxcirc \Cal{X}_k:= (\cdots (( \Cal{X}_1 \boxcirc \Cal{X}_2) \boxcirc \Cal{X}_3) \cdots)  \boxcirc \Cal{X}_k \] and obtaining grouping maps from the universal map in (\ref{eq:colim-circ}). Our notion of $\boxcirc$-monoids are \textit{normal oplax monoids} with respect to such structure by appealing to Ching \cite[3.4]{Ching-oplax}, noting in particular that four-fold and higher associativity diagrams are known to commute given the commutativity of (\ref{eq:assoc-boxcirc}).
\end{remark}

\begin{proposition}
    \label{def:box-circ-prod-map}
    The cosimplicial symmetric sequence $C(\Cal{O})$ (see \eqref{al:CO}) admits a natural $\boxcirc$-monoid structure, i.e., there are maps $m\colon C(\Cal{O})\boxcirc C(\Cal{O})\to C(\Cal{O})$ and $u \colon \underline{I}\to C(\Cal{O})$ which satisfy associativity and unitality. 
\end{proposition}

\begin{proof}
    The map $m$ is that constructed in Example \ref{ex:C(O)-box}. The unit $I\to J$ provides a coaugmentation $I\to C(\Cal{O})$ which in turn induces a map $u\colon \underline{I}\to C(\Cal{O})$. 
    
    Associativity (\ref{eq:assoc-boxcirc}) follows from a routine calculation, observing that \[d^0 \colon (\fco \boxcirc \fco)^{q} \to (\fco \boxcirc \fco)^{q+1}\] is induced by $d^0 \circ \id\colon \fco^r\circ \fco^s\to \fco^{r+1} \circ \fco^s$ for $r+s=q$. Similarly, the right-hand triangle from the unitality diagram (\ref{eq:unit-boxcirc}) is granted by the following commuting diagrams\[ \xymatrix{
        \underline{I}^p \circ C(\Cal{O})^q \ar[r]^-{u^p\circ \id} \ar[d]^-{\cong}
        & 
        C(\Cal{O})^p \circ \fco^q \ar[d]^-{m_{p,q}}
        \\ 
        \fco^q \ar[r]^-{d^0\cdots d^0} 
        & 
        \fco^{p+q}
        }\]
    for all $p,q$. A similar argument provides the commutativity of the other side of the unitality diagram. 
\end{proof}

Theorem \ref{main}(a) is then obtained as a corollary to the following proposition, the proof of which is deferred to Section \ref{sec:1a}. As such, the aim of the following sections is to set up a precise framework to describe what is meant by $A_{\infty}$-operad.

\begin{proposition}  \label{main-1}
    If $\Cal{X}$ is a $\boxcirc$-monoid in $\SymSeq_{\Spt}^{\bdel}$, then $\Tot \Cal{X}$ is an $A_{\infty}$-monoid with respect to the composition product (i.e., $A_{\infty}$-operad). 
\end{proposition}